\documentclass[11pt]{article}
\usepackage[margin=.9in]{geometry}
\usepackage{amsmath,amssymb,amsthm}
\usepackage[T1]{fontenc}
\usepackage{lmodern,microtype,hyperref}
\newtheorem{theorem}{Theorem}
\newtheorem{lemma}[theorem]{Lemma}
\newtheorem{proposition}[theorem]{Proposition}
\newtheorem{corollary}[theorem]{Corollary}
\theoremstyle{remark}\newtheorem{remark}[theorem]{Remark}
\newcommand{\norm}[1]{\left\lVert#1\right\rVert}
\DeclareMathOperator{\diag}{diag}\DeclareMathOperator{\ran}{ran}
\newcommand{\R}{\mathbb R}\newcommand{\C}{\mathbb C}
\begin{document}
\section{Strict support of the physical stationary background}
\label{st-section}

Fix odd $L$, dimension four, color and the compatible periodic source
geometry with distinct block interfaces. Current spacing is one, old
fine spacing is $M^{-1}$ with $M=L^k$, and next-coarse spacing is $L$.
Write $V=e^{ia}$ and $W=e^{iw}$. These $a,w$ are link angles. In current
physical units the next-coarse potential is $w/L$; in next-coarse unit
coordinates it is $w$. All formulas below use $w$, not $w/L$.
Constants may depend on fixed $L$, but not on $M$ or volume.

Group current coordinates into $L$-block cells, indexed by next-coarse
cells. The retained nonpivot coordinate $z$ and next-coarse datum $w$
have fixed real dimensions per block
\begin{equation}
 D_z=d(L^d-1)\dim\mathfrak g,\qquad D_w=d\dim\mathfrak g,
 \qquad D=D_z+D_w=dL^d\dim\mathfrak g.                         \label{st-dimensions}
\end{equation}
No fluctuation coupling has yet been inserted in $z$. Fixed-dimensional
norm equivalences are included in all radii. Scalar Cauchy estimates use
an actual coordinate product polydisc after this conversion.

We use the previously established ordinary source average and pivot
chart, the physical gauge-zero branch obtained by the two minimum
comparisons, and the strict scalar classical-action expansion. The flat
retained precision has the uniform bounds
\begin{equation}
 A_{M,0}=K_L^T(q_M+T_L^TT_L)K_L,\qquad C_{M,0}=A_{M,0}^{-1},
 \qquad \gamma_*I\le A_{M,0}\le B_*I,                         \label{st-flat-precision}
\end{equation}
and both matrices have uniformly exponentially summable block rows and
columns. Strict scalar support is available with any prescribed finite
cardinality reserve, at the cost of a smaller uniform product radius.
The result proved here is one uniform classical-background transport,
not a claim of closure of repeated quantum RG steps.

\subsection{A local affine pivot chart}
\label{st-affine-section}

At the identity background let
\begin{equation}
 F_1(a)=\frac1i\log M_L(e^{ia})=A_1a+C_1(a),\qquad
 C_1(0)=DC_1(0)=0.                                           \label{st-constraint}
\end{equation}
The actual source normalization is
\begin{equation}
 A_1=LQ_L^{\rm flat},\qquad A_1h_1=I,\qquad A_1K_L=0,
 \qquad h_1(c)=L^{d-1}I\text{ at }p(c).                       \label{st-pivot-normalization}
\end{equation}
The pivot coefficient in $Q_L^{\rm flat}$ is $L^{-d}$, hence its
coefficient in $A_1$ is $L^{1-d}$. Nonpivot restriction of $K_L$ is the
identity; $h_1$ has only pivot entries.

Prescribe the next-coarse datum $w$ and the retained absolute logarithmic
coordinates $z$ by solving
\begin{equation}
 a=K_Lz+h_1w-h_1C_1(a).                                     \label{st-affine-equation}
\end{equation}
Each ordinary constraint equation reads its own pivot and no other
pivot. Thus this is a family of independent finite-dimensional
Lie-algebra-block implicit equations, with neighboring nonpivot values
as parameters. Their derivative in the pivot at zero is invertible.
There are only finitely many local types. A common small contraction
therefore gives a holomorphic solution $a=\mathfrak a(w,z)$ satisfying
\begin{align}
 F_1(\mathfrak a(w,z))&=w,&
 \mathfrak a(w,z)|_{\rm nonpivot}&=z,\notag\\
 \|\mathfrak a(w,z)\|_\infty
    &\le C_L(\|w\|_\infty+\|z\|_\infty),&
 \mathfrak a(w,z)&=K_Lz+h_1w+O((\|w\|+\|z\|)^2).
                                                               \label{st-affine-bounds}
\end{align}
Iteration does not spread to another pivot equation. The dependency
neighborhood and all local analytic constants are fixed at $L$ and
independent of $M$ and volume. This is the unique small constraint
solution with the specified nonpivot values, equivalently the
identity-background ambient source straightening evaluated at
$K_Lz+h_1w$.

\subsection{Scalar support and the exact flat mixed term}
\label{st-scalar-section}

Let $G$ be the actual averaged-contour exponential gauge action, in the
coupling-free normalization, and define
\begin{equation}
 T_M(w,z)=S_M(\mathfrak a(w,z))+G(e^{i\mathfrak a(w,z)}).
                                                               \label{st-action}
\end{equation}
The strict scalar support theorem applies to this fixed-range analytic
substitution and the local gauge action. After fixed $L$-block
regrouping it yields
\begin{align}
 T_M(w,z)&=\sum_{X\ne\varnothing}T_X(w|_X,z|_X),             \label{st-scalar-expansion}\\
 \mathcal K&:=\sup_i\sum_{X\ni i}
 e^{\sigma\tau(X)+A_{
m c}|X|}\|T_X\|_\rho<\infty.         \label{st-scalar-norm}
\end{align}
Here $\tau$ is graph Steiner length on the $L$-block graph, allowing
auxiliary vertices, $\sigma>0$, and $A_{
m c}>0$ is any prescribed
finite reserve. The coordinate-polydisc radius $\rho$ and $\mathcal K$
are uniform in $M$ and volume. Every coefficient starts at degree at
least two: the action does so, the chart vanishes at zero, and the gauge
action has no constant or linear part at the identity.

Put $P_M=q_M+T_L^TT_L$. The full scalar action is stationary at zero, so
only the linear chart enters its quadratic term:
\begin{equation}
 T_M^{(2)}(w,z)=\tfrac12
 \langle K_Lz+h_1w,P_M(K_Lz+h_1w)\rangle.                    \label{st-quadratic-action}
\end{equation}
Thus
\begin{align}
 D_zT_M(w,z)&=A_{M,0}z+B_Mw+N_M(w,z),\notag\\
 B_M&=K_L^TP_Mh_1=K_L^Tq_Mh_1,&
 N_M(0,0)&=DN_M(0,0)=0.                                    \label{st-gradient-split}
\end{align}
The last simplification follows from $T_Lh_1=0$, because the root
contours are block internal and contain no pivot. The rectangular
$B_M$ has uniform exponential block bounds. The pairings here are
counting pairings of the stated angle coordinates; no additional
$L^4$ transpose factor is introduced. Replacing $w$ by the next-coarse
physical potential means substituting $w=L a_{\rm next}$ explicitly.

\subsection{An analytic mixed-slot bound from scalar supports}
\label{st-gradient-section}

The supports in \eqref{st-scalar-expansion} are not uniformly bounded.
The finite-range vertex lemma therefore cannot be applied directly to
their gradient. In particular, crude polarization followed by
$|X|^n$ slot counting can produce a factorial loss. Coordinatewise
Cauchy gives the needed stronger conclusion.

Write $u=(w,z)$ in a fixed scalar basis. On the coordinate polydisc
there are $D|X|$ variables. If $M_X=\|T_X\|_\rho$, its scalar Taylor
coefficients satisfy
\begin{equation}
 |c_{X,\beta}|\le M_X\rho^{-|\beta|}.                       \label{st-scalar-cauchy}
\end{equation}
For one retained derivative coordinate $\alpha$ at a cell $i\in X$,
differentiating the product of geometric series gives, for $t<\rho$,
\begin{equation}
 \sum_{\beta\ge e_\alpha}
 \beta_\alpha |c_{X,\beta}|t^{|\beta|-1}
 \le\frac{M_X}{\rho}(1-t/\rho)^{-D|X|-1}.                   \label{st-gradient-cauchy}
\end{equation}
Sum the at most $D_z$ output components. A symmetric multilinear Taylor
coefficient distributes each scalar monomial coefficient among its
ordered slots with multinomial weights summing to one. Consequently
the sum of its ordered cell-slot operator norms is bounded by the same
absolute scalar coefficient sum, with fixed norm-conversion constants
absorbed in this bound. Every output and input-slot root lies in $X$,
so their rooted Steiner length is at most $\tau(X)$.

For $0<\kappa\le\sigma$, the supremum of the all-order positive row
series is therefore bounded by
\begin{equation}
 \mathcal G(t)=\frac{D_z}{\rho}\sup_i\sum_{X\ni i}
 e^{\kappa\tau(X)}M_X(1-t/\rho)^{-D|X|-1}.                  \label{st-row-series}
\end{equation}
Choose $t_0<\rho/2$ such that
\begin{equation}
 D\log\frac1{1-t_0/\rho}\le\frac{A_{
m c}}2.
 \qquad\text{Then}\qquad
 \mathcal G(t_0)\le\frac{2D_z\mathcal K}{\rho}=:M_g.         \label{st-majorant-radius}
\end{equation}
The supremum over output cells cannot be interchanged with the order
sum. Instead let $v_n$ be the supremum of the order-$n$ mixed-slot row
norm. Equation \eqref{st-row-series} gives the separate bound
$v_n\le M_gt_0^{-n}$. Removing degrees zero and one gives the positive
majorant for $N_M$:
\begin{equation}
 \Phi(t)=\sum_{n\ge2}v_nt^n
 \le\frac{M_g(t/t_0)^2}{1-t/t_0},\qquad
 \Phi'(t)\le\frac{8M_gt}{t_0^2}\quad(0\le t\le t_0/2).
                                                               \label{st-nonlinear-majorant}
\end{equation}
This deliberate geometric factor makes the order/root-supremum
bookkeeping valid. It leaves an analytic quadratic majorant uniform in
refinement and volume.

The original $X$ is a fixed coefficient carrier at the flat baseline,
not an additional variable support. After substituting families for
$w,z$, their union is the actual variable support. Internal rooted
connections through $X$ have already been paid by
\eqref{st-row-series}. This preserves strict dependence on the original
next-coarse variables.

\subsection{The support-valued stationary fixed point}
\label{st-contraction-section}

For a prescribed finite $a_{
m c}\ge0$, use
\begin{equation}
 \|u\|_{\kappa,a_{
m c};r}
 =\sup_i\sum_X e^{\kappa\tau(\{i\}\cup X)+a_{
m c}|X|}
             \sup_{\|w|_X\|<r}\|u_{i;X}(w|_X)\|.            \label{st-rooted-norm}
\end{equation}
The source $w$ has singleton norm $e^{a_{
m c}}r$. For joint insertion
in $N_M$, use the maximum of the separate $w$- and $z$-family norms and
color slots by scalar coordinate and field type. The positive scalar
majorant above already sums all such colors and types, while each
scalar-component family is bounded by its vector-family norm. This
does not identify that maximum with
$\sum_X\|(w_X,z_X)\|$, which could be larger.

Linear exponential block bounds act in \eqref{st-rooted-norm} by their
weighted row sums. The rooted Steiner union inequality and
$e^{a_{
m c}|\cup X_r|}\le\prod_r e^{a_{
m c}|X_r|}$ allow the
proved support-substitution lemma to use
\eqref{st-nonlinear-majorant}. Stationarity becomes
\begin{equation}
 z=-C_{M,0}B_Mw-C_{M,0}N_M(w,z).                            \label{st-stationary-fixed-point}
\end{equation}
Let $c_0$ bound the weighted norm of $C_{M,0}$ and choose
$\ell_0\ge\max\{1,\|C_{M,0}B_M\|_\kappa\}$ uniformly. Put
\begin{equation}
 r_*\le\min\left\{\frac{t_0}{2},
                  \frac{t_0^2}{16c_0M_g}\right\},\qquad
 r\le\frac{r_*}{2\ell_0e^{a_{
m c}}}.                     \label{st-data-radius}
\end{equation}
Decrease these radii further so that reconstruction and evaluation stay
within the common prior chart and uniqueness neighborhoods. On the
family ball $\|z\|\le2\ell_0e^{a_{
m c}}r$, the map
\eqref{st-stationary-fixed-point} is a contraction with constant at most
$1/2$ and preserves that ball. Hence it has a unique holomorphic
supported solution with
\begin{equation}
 \|z_*\|_{\kappa,a_{
m c};r}\le2\ell_0e^{a_{
m c}}r,
 \qquad
 \|z_*+C_{M,0}B_Mw\|_{\kappa,a_{
m c};r}\le C_{a_{
m c}}r^2.
                                                               \label{st-retained-support}
\end{equation}
Evaluation gives a genuine stationary solution on each finite
regulator. Picard iteration defines an absolutely convergent rooted
expansion with original $w$-cell leaves.

The finite-range chart $\mathfrak a$ can now be substituted by the
finite-range support lemma. Define
\begin{equation}
 a_*(w)=\mathfrak a(w,z_*(w)),\qquad
 H_{{\rm cur},M}=h_1-K_LC_{M,0}B_M.                          \label{st-current-extension}
\end{equation}
Then
\begin{equation}
 \|a_*\|_{\kappa,a_{
m c};r}\le C_{a_{
m c}}r,
 \qquad
 \|a_*-H_{{\rm cur},M}w\|_{\kappa,a_{
m c};r}
                                      \le C_{a_{
m c}}r^2.          \label{st-current-support}
\end{equation}
This $H_{{\rm cur},M}$ is the current-angle extension, distinct from the
earlier original-fine extension. It has a uniform exponential
rectangular block bound and $A_1H_{{\rm cur},M}=I$. Its source gauge
condition will follow by identifying the physical branch below; no
gauge-zero conclusion is inferred from stationarity alone. Local
exponentiation also supplies strict families for $e^{ia_*}-I$ and its
inverse minus the identity. For activity transport the bound
$s(r)=C_{a_{
m c}}r$ on the potential itself is sufficient.

\subsection{Identification of the physical gauge-zero branch}
\label{st-identification-section}

For real small $w$, the already proved two-scale minimum comparison and
rooted averaged-contour gauge construction give a branch
$V_{{\rm old},*}(e^{iw})=e^{ia_{{\rm old},*}}$ with
\begin{align}
 M_L(V_{{\rm old},*})&=e^{iw},& P_{V_{{\rm old},*}}(y,x)&=1,\notag\\
 \|a_{{\rm old},*}\|_\infty&\le C_L\|w\|_\infty,&
 S_M(a)&\ge S_M(a_{{\rm old},*})                           \label{st-known-minimum}
\end{align}
for every real feasible $a$ in the common small comparison ball. The
actual gauge action is nonnegative on real unitary fields and vanishes
at $P=1$. Hence
\begin{equation}
 S_M(a)+G(e^{ia})\ge S_M(a_{{\rm old},*})+G(V_{{\rm old},*})
                                                               \label{st-action-comparison}
\end{equation}
on that feasible neighborhood. The affine chart parametrizes every
sufficiently small feasible field by its retained absolute coordinates.
Thus $z_{{\rm old},*}=a_{{\rm old},*}|_{\rm nonpivot}$ is a critical
point of $T_M(w,\cdot)$.

The scalar support expansion and Cauchy with fixed margins give
\begin{equation}
 \|D_z^2T_M(w,z)-A_{M,0}\|_\kappa
                \le C_L(\|w\|_\infty+\|z\|_\infty).         \label{st-convexity}
\end{equation}
For each coefficient, contract its third derivative with the bounded
direction $(w,z)$; summing the second retained matrix index costs only
$|X|$, paid by a cardinality spare. Its input/output distance is at
most $\tau(X)$. Symmetry gives the column bound. This calculation
differentiates the full scalar action, so both chart curvature and
actual gauge curvature remain included. It does not use the ordinary
rooted average's physical adjoint.

Choose a common real retained supremum ball and a data ball making
\eqref{st-convexity} at most $\gamma_*/2$. The retained Hessian is
positive on this convex ball, so there is at most one critical point.
Both $z_*$ from \eqref{st-stationary-fixed-point} and
$z_{{\rm old},*}$ lie inside after a uniform shrink. Therefore
\begin{equation}
 e^{ia_*(w)}=V_{{\rm old},*}(e^{iw}),\qquad
 M_L(e^{ia_*(w)})=e^{iw},\qquad P_{e^{ia_*(w)}}(y,x)=1.       \label{st-physical-identities}
\end{equation}
This identifies with a known gauge-zero minimum; a generic critical
equation was not used to assert the gauge condition. Holomorphy and the
finite-dimensional identity theorem extend these identities from the
real ball to the common complex product domain. Real minimum
inequalities are not extended as inequalities. Linearizing the last
identity gives $T_LH_{{\rm cur},M}=0$.

\subsection{Transport conclusion and its domain limitation}
\label{st-scope-section}

For any prescribed finite cardinality reserve and sufficiently small
positive $\kappa$, \eqref{st-current-support} is now a strict-support
expansion of the actual physical stationary background in its original
next-coarse angle variables. The radius is uniform in old refinement
and volume. Output roots are $L$-blocks of current fields and input
labels are next-coarse cells on the same block graph. If an activity
uses another current-cell convention, its projection to this graph and
the finite fiber multiplicity must be included explicitly; fixed
regrouping is not a contraction estimate. Coefficients can depend on
the fixed baseline and regulator. Strict variable support does not by
itself assert volume consistency of a chosen decomposition.

There is a concrete obstruction to inferring an $L^{-1}$ contraction
of the angle supremum domain. Take small coarse pure-gauge data
$W(y,y')=g_yg_{y'}^{-1}$ and set the current-site gauge to $g_y$
throughout each $L$-block. Its current field is the identity internally
and equals $W$ on the corresponding block-boundary links. It has
$P_V=1$, ordinary average $W$ and zero classical action. By the same
minimum comparison and uniqueness it is the physical stationary
representative. Its boundary angle is of order $\|w\|$, not
$\|w\|/L$. The physical conversion $a_{\rm next}=w/L$ does not change
that fact about the angle-domain map.

This theorem supplies one uniform step and an explicit small-domain
bound $s(r)=C_ar$ for substitution in supplied analytic activities.
Preservation of a quantum activity norm, any associated radius loss,
and a compatible common domain under repeated integrations require
separate proofs. No running-coupling, mixed-history, large-field or
interacting continuum conclusion follows here.

\end{document}
